\documentclass[11pt]{article}
\usepackage[margin=1in]{geometry}
\usepackage{amsmath,amssymb,amsthm}
\usepackage{booktabs}
\usepackage[T1]{fontenc}
\usepackage[utf8]{inputenc}
\title{D5 note: T3 repeated-end corner via four explicit collision families}
\author{}
\date{2026-03-21}

\newtheorem{theorem}{Theorem}
\newtheorem{proposition}[theorem]{Proposition}
\newtheorem{corollary}[theorem]{Corollary}
\newtheorem{remark}[theorem]{Remark}

\newcommand{\eX}{e_0}
\newcommand{\eQ}{e_1}
\newcommand{\eW}{e_2}
\newcommand{\eV}{e_3}
\newcommand{\eU}{e_4}

\begin{document}
\maketitle

\section*{Purpose}

For the repeated-end corner
\[
T3:\qquad \sigma_{33}=[4,4,1]/[2,2,1],
\]
this note extracts the local collision lemma package directly from the surfaced
2026-03-21 bundle.
The point is to replace the earlier dossier statement
``the checked collision families look clean'' by an actual theorem-facing proof.
The output is twofold:
\begin{enumerate}
\item an explicit four-family collision theorem for the base row, valid for every
odd modulus $m\ge 5$;
\item an immediate corollary that the entire \emph{one-label} tiny-repair scope of
artifact \texttt{032} (one-gate and one-bit repairs, with or without the bundled
cocycle-defect toggles) is already impossible for this row.
\end{enumerate}

This does \emph{not} claim anything about richer multi-class repair transducers.
It only internalizes the one-gate / one-bit slice of the tiny-repair background
for the specific row $\sigma_{33}$.

\section*{Coordinate conventions and surfaced background rule}

Write a color-$0$ vertex as
\[
x=(x_0,q,w,v,u)\in (\mathbb Z/m\mathbb Z)^5,
\]
and set
\[
\ell(x)=x_0+q+w+v+u \pmod m,
\qquad
s(x)=w+u \pmod m.
\]
Let $\eX,\eQ,\eW,\eV,\eU$ denote the standard coordinate increments in the
$(x_0,q,w,v,u)$ coordinates.
Thus the surfaced bundle step maps are simply
\[
x\mapsto x+\eX,\ x+\eQ,\ x+\eW,\ x+\eV,\ x+\eU.
\]

The surfaced mixed background rule comes from
\texttt{torus\_nd\_d5\_endpoint\_latin\_repair.py} via
\texttt{\_mixed\_rule()}, which reads the simple witness with
\begin{itemize}
\item layer-$2$ table $(0,2)$,
\item layer-$3$ mode tables $(0,3)$ and $(3,0)$,
\item predecessor flag \texttt{pred\_sig1\_wu2}.
\end{itemize}
For color $0$, this means:
\begin{enumerate}
\item every layer-$1$ vertex moves by $\eU$;
\item every layer-$2$ vertex with $q=m-1$ moves by $\eW$;
\item every layer-$3$ vertex with $s\neq 2$ moves by
\[
\begin{cases}
\eV, & q=m-1,\\
\eX, & q\neq m-1.
\end{cases}
\]
\end{enumerate}
We will only use these slices, so the $s=2$ layer-$3$ cases are irrelevant here.
Call this background row $F$.

\section*{The $\sigma_{33}$ row and its six labelled source classes}

Fix the representative seed $(w_0,s_0)=(0,0)$.
For $m\ge 5$ define the two layer-$1$ endpoint families
\[
A:=\{x:\ \ell(x)=1,\ q=m-1,\ w=m-1,\ s(x)\neq 0\},
\]
\[
B:=\{x:\ \ell(x)=1,\ q=m-1,\ w=0,\ s(x)\neq 0\}.
\]
The color-$0$ labelled source classes for
\[
\sigma_{33}=[4,4,1]/[2,2,1]
\]
are then exactly
\[
L1=A,
\qquad
R1=B,
\qquad
L2=A+\eU,
\qquad
R2=B+\eW,
\qquad
L3=A+2\eU,
\qquad
R3=B+2\eW.
\]
Equivalently, the six labelled classes are the following affine slices:
\[
\begin{array}{rcl}
L1&=&\{\ell=1,\ q=m-1,\ w=m-1,\ s\neq 0\},\\
L2&=&\{\ell=2,\ q=m-1,\ w=m-1,\ s\neq 1\},\\
L3&=&\{\ell=3,\ q=m-1,\ w=m-1,\ s\neq 2\},\\
R1&=&\{\ell=1,\ q=m-1,\ w=0,\ s\neq 0\},\\
R2&=&\{\ell=2,\ q=m-1,\ w=1,\ s\neq 1\},\\
R3&=&\{\ell=3,\ q=m-1,\ w=2,\ s\neq 2\}.
\end{array}
\]
Each family has cardinality $m(m-1)$.
Indeed, in $A$ and $B$ one may choose $(x_0,v)$ arbitrarily, then $u$ is forced by
$\ell=1$, and the single forbidden value of $u$ coming from $s\neq 0$ excludes
exactly one diagonal of size $m$ among the $m^2$ choices.
Translation by $\eU$ or $\eW$ preserves cardinality.

Let $G$ denote the color-$0$ row defined by $\sigma_{33}$.
From the surfaced source-builder one reads immediately:
\[
G=F\text{ on }L1\cup R2,
\]
and on the remaining four classes
\[
G(x)=x+\eW\text{ on }R1,
\qquad
G(x)=x+\eU\text{ on }L2,
\qquad
G(x)=x+\eQ\text{ on }L3\cup R3.
\]
So the repeated-end corner differs from the background row only on
\[
R1,\ L2,\ L3,\ R3.
\]

\section*{The four explicit collision families}

\begin{theorem}[four explicit color-$0$ collision families for $T3$]
\label{thm:T3-four-families}
Let $m\ge 5$ be odd.
For the color-$0$ row $G$ of $\sigma_{33}=[4,4,1]/[2,2,1]$ at $(w_0,s_0)=(0,0)$,
the only colliding changed labels are
\[
R1,\ L2,\ L3,\ R3,
\]
and each contributes exactly $m(m-1)$ background-vs-changed collisions.
More precisely:
\begin{center}
\begin{tabular}{@{}llll@{}}
\toprule
changed source $x$ & background mate $y$ & common target $t$ & collision type \\
\midrule
$x\in R1$ & $y=x+\eW-\eU$ & $t=x+\eW$ & $(B,R1)$ \\
$x\in L2$ & $y=x-\eW+\eU$ & $t=x+\eU$ & $(B,L2)$ \\
$x\in L3$ & $y=x-\eX+\eQ$ & $t=x+\eQ$ & $(B,L3)$ \\
$x\in R3$ & $y=x-\eX+\eQ$ & $t=x+\eQ$ & $(B,R3)$ \\
\bottomrule
\end{tabular}
\end{center}
For every row in the table,
\[
G(x)=F(y)=t,
\]
with $y$ a genuine background source.
The four target sets are pairwise disjoint, so color $0$ has exactly
\[
4m(m-1)
\]
overfull targets, namely $m(m-1)$ of each pair type
\[
(B,R1),\ (B,L2),\ (B,L3),\ (B,R3).
\]
By color rotation, the full five-color row has exactly
\[
20m(m-1)
\]
overfull targets.
\end{theorem}

\begin{proof}
We treat the four families one by one.

\smallskip
\noindent\textbf{R1.}
Take $x\in R1$.
Then $x$ has layer $1$, $q=m-1$, $w=0$, and $s(x)\neq 0$.
Since $G$ uses the right first-step direction $2$, we have
\[
G(x)=x+\eW.
\]
Set $y=x+\eW-\eU$.
Then $y$ is still layer $1$, still has $q=m-1$, now has $w=1$, and hence is not in
any changed source class.
Because every layer-$1$ background source moves by $\eU$,
\[
F(y)=y+\eU=x+\eW=G(x).
\]
The target $t=x+\eW$ lies in the $R2$ slice $(\ell=2,q=m-1,w=1)$.
Different $x\in R1$ give different $t$, so this produces exactly $|R1|=m(m-1)$
collisions of type $(B,R1)$.

\smallskip
\noindent\textbf{L2.}
Take $x\in L2$.
Then $x$ has layer $2$, $q=m-1$, $w=m-1$, and $s(x)\neq 1$.
Since $G$ uses the left second-step direction $4$, we have
\[
G(x)=x+\eU.
\]
Set $y=x-\eW+\eU$.
Then $y$ is still layer $2$, still has $q=m-1$, now has $w=m-2$, and therefore is
not in $L2$ or $R2$.
For layer $2$ with $q=m-1$ the background row uses $\eW$, so
\[
F(y)=y+\eW=x+\eU=G(x).
\]
The target $t=x+\eU$ lies in the $L3$ slice $(\ell=3,q=m-1,w=m-1)$.
Again the map $x\mapsto t$ is injective, so we get exactly $|L2|=m(m-1)$
collisions of type $(B,L2)$.

\smallskip
\noindent\textbf{L3.}
Take $x\in L3$.
Then $x$ has layer $3$, $q=m-1$, $w=m-1$, and $s(x)\neq 2$.
Since $G$ uses the third-step direction $1$, we have
\[
G(x)=x+\eQ.
\]
Set $y=x-\eX+\eQ$.
Then $y$ is still layer $3$, has $q=0$, keeps the same $w$ and the same
$s\neq 2$, so $y$ is not in any changed source class.
On layer $3$ with $q\neq m-1$ and $s\neq 2$, the background row uses $\eX$.
Hence
\[
F(y)=y+\eX=x+\eQ=G(x).
\]
The target $t=x+\eQ$ lies in layer $4$ with $(q,w)=(0,m-1)$.
Therefore we obtain exactly $|L3|=m(m-1)$ collisions of type $(B,L3)$.

\smallskip
\noindent\textbf{R3.}
Take $x\in R3$.
Then $x$ has layer $3$, $q=m-1$, $w=2$, and $s(x)\neq 2$.
Again $G(x)=x+\eQ$.
Set $y=x-\eX+\eQ$.
Then $y$ is still layer $3$, now has $q=0$, keeps the same $w=2$ and the same
$s\neq 2$, so it is background.
As above,
\[
F(y)=y+\eX=x+\eQ=G(x).
\]
The target now lies in layer $4$ with $(q,w)=(0,2)$, so this family is disjoint
from the $(B,L3)$ family as soon as $m\ge 5$.
Thus we obtain exactly $|R3|=m(m-1)$ collisions of type $(B,R3)$.

\smallskip
We have now constructed $4m(m-1)$ pairwise disjoint targets carrying the stated
collision types.
Since the background row $F$ is Latin, every overfull target of $G$ must involve
at least one changed source, and the row differs from $F$ only on the four changed
classes $R1,L2,L3,R3$.
Moreover the changed-image sets found above are pairwise disjoint, so no target can
contain two changed sources from different families.
Therefore the listed collisions exhaust all color-$0$ overfull targets.
Color rotation gives the same count on every color.
\end{proof}

\begin{corollary}[row-level obstruction for the repeated-end corner]
For every odd $m\ge 5$, the repeated-end row
\[
\sigma_{33}=[4,4,1]/[2,2,1]
\]
is non-Latin already at the base row.
\end{corollary}

\begin{proof}
Theorem~\ref{thm:T3-four-families} gives explicit overfull targets.
\end{proof}

\section*{One-gate and one-bit repairs are already impossible}

Artifact \texttt{032} allows two one-label repair modes on top of the labelled row:
\begin{enumerate}
\item a one-gate repair, which reroutes every source in one chosen label class;
\item a one-bit repair, which reroutes a selected bit-slice inside one chosen label
class.
\end{enumerate}
It also allows optional cocycle-defect toggles on two omitted layer-$3$ slices.
For $(w_0,s_0)=(0,0)$ these toggles touch only the slices
\[
O_L=\{\ell=3,\ q=m-1,\ w=0,\ s=2\},
\qquad
O_R=\{\ell=3,\ q=m-1,\ w=1,\ s=2\}.
\]

\begin{proposition}[the whole one-label repair scope fails for $T3$]
\label{prop:T3-one-label-fails}
Let $m\ge 5$ be odd.
Take the repeated-end row $\sigma_{33}=[4,4,1]/[2,2,1]$ and apply any one-gate or
one-bit repair from artifact \texttt{032}, with cocycle defect chosen from
\texttt{none}, \texttt{left}, \texttt{right}, or \texttt{both}.
The resulting row is still non-Latin.
\end{proposition}

\begin{proof}
A one-gate or one-bit repair changes sources inside at most one label class.
The cocycle-defect toggles change only the omitted slices $O_L$ and $O_R$.
But the four collision families from
Theorem~\ref{thm:T3-four-families} use labels
\[
R1,\ L2,\ L3,\ R3,
\]
with source slices
\[
R1:\ (\ell,w)=(1,0),
\qquad
L2:\ (\ell,w)=(2,m-1),
\qquad
L3:\ (\ell,w,s)=(3,m-1,\neq 2),
\qquad
R3:\ (\ell,w,s)=(3,2,\neq 2),
\]
and background mates lying in the slices
\[
(\ell,w)=(1,1),
\qquad
(\ell,w)=(2,m-2),
\qquad
(\ell,q,w,s)=(3,0,m-1,\neq 2),
\qquad
(\ell,q,w,s)=(3,0,2,\neq 2).
\]
None of these sources or background mates lies in $O_L$ or $O_R$.
Therefore the cocycle-defect toggles do not touch them.
A one-label repair can at most destroy the collision family attached to the chosen
label, and it leaves the other three families unchanged.
Hence after any such repair at least three explicit collision families survive, so
the row remains non-Latin.
\end{proof}

\begin{remark}[scope]
This proposition matches the negative one-gate / one-bit scan from the bundle, but
it does \emph{not} eliminate richer multi-class repair mechanisms.
So it should be read as:
\begin{quote}
for $T3$, the entire one-label tiny-repair branch is now theorem-level dead.
\end{quote}
If the ambient tiny-repair background $T0$ is later unpacked exactly as the same
one-label repair scope, then $T3$ is fully closed.
If $T0$ allows a richer repair class, then the present note closes only the one-label
slice of that background.
\end{remark}

\begin{remark}[the exceptional modulus $m=3$]
The proof is stated for $m\ge 5$ because at $m=3$ the slices with $w=2$ and
$w=m-1$ coincide, and the clean four-family decomposition degenerates.
The bundle computation does indeed show a different color-$0$ profile at $m=3$.
Nothing in this note uses $m=3$.
\end{remark}

\end{document}
